% Generated draft: requires mathematical and source review.
\documentclass[11pt]{article}
\usepackage[margin=25mm]{geometry}
\usepackage{fontspec}
\setmainfont{Latin Modern Roman}
\usepackage{amsmath,amssymb,mathtools,unicode-math}
\setmathfont{Latin Modern Math}
\usepackage{enumitem,needspace,xurl,hyperref}
\setlength{\parindent}{0pt}
\setlength{\parskip}{5pt}
\setlength{\emergencystretch}{3em}
\newcommand{\N}{\mathbb{N}}
\newcommand{\Z}{\mathbb{Z}}
\newcommand{\Q}{\mathbb{Q}}
\newcommand{\R}{\mathbb{R}}
\newcommand{\C}{\mathbb{C}}
\newcommand{\F}{\mathbb{F}}
\newcommand{\A}{\mathbb{A}}
\newcommand{\PP}{\mathbb{P}}
\newcommand{\E}{\mathbb{E}}
\newcommand{\eps}{\varepsilon}
\newcommand{\dd}{\,\mathrm{d}}
\begin{document}
\section*{UnsolvedMath ID: 30003643}
\textbf{Source problem code:} catalogue entry
\subsection{Definitions and mathematical statement}
For a positive integer \(n\), let \(\mathcal S(\R^n)\) denote the Schwartz class of smooth functions \(f:\R^n\to\C\) such that
\[
\sup_{x\in\R^n}|x^\alpha\partial^\beta f(x)|<\infty
\qquad\text{for all }\alpha,\beta\in\Z_{\geq0}^n.
\]
Here \(|x|\) is the Euclidean norm. For definiteness in this draft, write
\[
\widehat f(\xi)=\int_{\R^n}f(x)e^{-2\pi i x\cdot\xi}\,\dd x.
\]
The source does not specify the Fourier normalization, the scalar field, or the range of dimensions intended; the conventions above are explicit drafting conventions requiring review.

The supplied question asks whether, for every pair of sequences of nonnegative radii \((r_j)_{j\geq1}\) and \((\rho_j)_{j\geq1}\) satisfying
\[
\lim_{j\to\infty}r_j=\lim_{j\to\infty}\rho_j=+\infty,
\]
and every \(f\in\mathcal S(\R^n)\), the conditions
\[
\begin{aligned}
&\forall j\geq1\ \forall x\in\R^n,
&&|x|=r_j\ \Longrightarrow\ f(x)=0,\\
&\forall j\geq1\ \forall \xi\in\R^n,
&&|\xi|=\rho_j\ \Longrightarrow\ \widehat f(\xi)=0
\end{aligned}
\]
imply
\[
\forall x\in\R^n,\qquad f(x)=0.
\]
The displayed formulation records the literal universal reading of the supplied question; it does not establish that this reading is the intended or mathematically valid uniqueness assertion. The background describes the sequences as increasing, whereas the original and clean statements require only divergence to infinity. No spacing, density, coupling, or additional restrictions on the radii are supplied. It therefore remains unresolved whether the intended target is this unrestricted implication or a uniqueness question under omitted hypotheses. The source also does not state whether radius zero is allowed.
\subsection{Short English statement}
The supplied problem asks whether a Schwartz function must vanish identically when both it and its Fourier transform vanish on prescribed sequences of entire spheres whose radii tend to infinity. The intended dimensions and any necessary restrictions on the radii require clarification. The supplied open status is dated curation, not an independently verified mathematical conclusion.
\subsection{Sources}
\begin{itemize}
\item Supplied local UnsolvedMath registry record, permanent ID \(30003643\), title ``Fourier Uniqueness from Spherical Zero Sets,'' collection ``Oberwolfach Reports Open Problems,'' exact source problem code \texttt{catalogue entry}. The registry gives proposed year 2017 and no named proposer; its literature-source label associates the problem with Hedenmalm.
\item Supplied source citation: \emph{Mini-Workshop: Interplay between Number Theory and Analysis for Dirichlet Series (2018)}, Oberwolfach Reports, pp. 3035--3069, DOI \texttt{10.4171/owr/2017/51}. \url{https://doi.org/10.4171/owr/2017/51}. The background describes Oberwolfach Reports as documenting mathematical workshops and their open-problem discussions. The supplied literature list labels this reference ``Oberwolfach Report 51, Hedenmalm problem.'' The citation's 2018 date and the registry's 2017 proposed year are preserved without resolving their relationship.
\item Supplied literature reference: Ramos--Stoller, \emph{Perturbed Fourier uniqueness in higher dimensions}. \url{https://arxiv.org/abs/2103.12015}.
\item Supplied literature reference: Hedenmalm, \emph{Deep zero problems}. \url{https://doi.org/10.1007/s11854-025-0376-1}.
\item Literature assessment dated exactly \texttt{2026-08-21}: the supplied curation labels the problem open and reports that the simultaneous Fourier-uniqueness question for arbitrary increasing radius sequences remains open, with published results covering structured special cases such as radii near square roots of integers. This is supplied dated curation derived from the supplied OWR curation, not independently verified here; both the status and the cited claims require independent verification.
\item The registry labels the statement \texttt{corrected\_verified} and reports manual semantic, typographic, and LaTeX review. These are supplied review claims and do not establish mathematical correctness or resolve the formulation issues recorded here.
\end{itemize}
\end{document}
